% components/frontmatter.tex — Imprint, preface, TOC
% ----------------------------------------------------
% components/cover.tex — Title page
% -----------------------------------
% Minimal academic cover: series band, title, rule, subtitle, author,
% institution, edition badge. Swap the TikZ motif to re-theme.
\begin{titlepage}
  \thispagestyle{empty}
  \centering
  \vspace*{18mm}

  % Series band
  {\color{navy}\rule{0.62\textwidth}{0.5pt}}\\[6pt]
  {\sffamily\small\bfseries\color{slate}\MakeUppercase{\CurriculumSeries}}\\[3pt]
  {\color{navy}\rule{0.62\textwidth}{0.5pt}}

  \vspace{16mm}

  % Geometric motif: parabola envelope + secant (inequalities theme)
  \begin{tikzpicture}[scale=1.0, line width=0.9pt]
    \draw[color=navy,   domain=-1.9:1.9, smooth] plot (\x, {0.42*\x*\x});
    \draw[color=gold]   (-1.6, {0.42*1.6*1.6}) -- (1.6, {0.42*1.6*1.6});
    \fill[color=gold]   (0, {0.42*1.6*1.6}) circle (1.6pt);
    \draw[color=slate,  line width=0.5pt] (-2.25,0) -- (2.25,0);
  \end{tikzpicture}

  \vspace{14mm}

  % Title block
  {\fontsize{30}{34}\selectfont\sffamily\bfseries\color{navy}
    \CurriculumTitle\par}
  \vspace{7mm}
  {\large\itshape\color{slate}\CurriculumSubtitle\par}

  \vspace{18mm}

  % Author & institution
  {\large\CurriculumAuthor\par}
  \vspace{2mm}
  {\normalsize\color{slate}\CurriculumInstitution\par}

  \vfill

  % Edition badge + date
  {\sffamily\small\bfseries\color{white}%
    \tcbox[on line, boxrule=0pt, arc=2pt, left=8pt, right=8pt, top=3.5pt,
           bottom=3.5pt, colback=navy, colframe=navy]
      {\CurriculumEdition}}\par
  \vspace{5mm}
  {\small\color{slate}\CurriculumDate\quad\textbullet\quad\CurriculumVersion\par}

  \vspace{10mm}
\end{titlepage}


% Imprint / edition notes page (verso of cover).
\cleardoublepage
\thispagestyle{empty}
\noindent{\sffamily\bfseries\color{navy}\CurriculumTitle}\\[2pt]
{\small\color{slate}\CurriculumSubtitle\par}
\smallskip
\noindent{\color{rulegray}\rule{\textwidth}{0.4pt}}\par\smallskip
{\small
\noindent\textbf{Edition.} \CurriculumEdition{} (\CurriculumVersion),
compiled \CurriculumDate.\par
\smallskip
\noindent\textbf{About this series.} This volume belongs to the
\CurriculumSeries{}, a structured program of study for competition
mathematics. Each chapter develops theory with complete proofs, worked
examples, and graded problem sets with hints and full solutions.\par
\smallskip
\noindent\textbf{Difficulty scale.}%
\begin{itemize}[leftmargin=1.6em, itemsep=2pt, topsep=3pt]
  \item \difficulty{1}\; Routine reinforcement of a definition or theorem.
  \item \difficulty{2}\; Standard olympiad application; one main idea.
  \item \difficulty{3}\; Hard: requires combining two or more ideas.
  \item \difficulty{4}\; Very hard; shortlist level.
  \item \difficulty{5}\; Olympiad summit (e.g.\ IMO Problem 6 caliber).
\end{itemize}\par}

\frontmatter
\cleardoublepage
\tableofcontents

\cleardoublepage
\section*{Preface}
\markright{Preface}
\noindent This chapter forms the opening unit of the inequalities
strand. It may be read independently; the only background assumed is
fluency with basic algebra and the willingness to check every step of
every proof.

\bigskip
\noindent\textit{How to use this book.} Read each proof actively: cover
it, attempt it, compare. Attempt every exercise before reading its
hint; hints are ordered so that earlier hints spoil less. Full
solutions appear at the end of each chapter.
